\section{Related Work}
\label{sec: related_work}

\noindent\textbf{Memory for LLM Agents.}
Memory has emerged as an essential module in modern agent systems~\citep{zhang2024survey, hu2025memory, he2026memoryarenabenchmarkingagentmemory} to enhance their performance by utilizing past information. Existing memory systems organize and store information in various forms, including plain text~\citep{packer2023memgpt}, latent embeddings~\citep{wang2025m}, and structured graphs~\citep{xu2025mem, chhikara2025mem0, li2025memos}. Beyond memory content, those methods usually involve retrieval mechanisms (e.g., semantic search) with memory management strategies (e.g., updating)~\citep{tan-etal-2025-prospect, hu-etal-2025-hiagent}. More recently, with the growing development of reinforcement learning (RL) in LLM agents, RL has also been leveraged for memory management in agent systems~\citep{yu2025memagent, zhou2025mem1}.
While most efforts primarily emphasize personalization~\citep{zhang2025prime, zhong2024memorybank} and long-context management~\citep{maharana-etal-2024-evaluating, wu2025longmemeval, hu2025evaluating}, 
this paper falls in the research line of learning from past experiences as memory~\citep{zhao2024expel, su2025learnbyinteract}, which is a critical aspect for developing self-evolving agent systems~\citep{gao2025survey, liang2024self}. Different from previous works that emphasize reusing successful trajectories~\citep{zheng2024synapse, tang2025chemagent}, procedural workflows~\citep{wang2025agent, qian2024investigate, fang2025memp, liu-etal-2025-contextual, 10.5555/3692070.3693628}, or instance-level concepts~\citep{10.1145/3690624.3709171, suzgun2025dynamic, zhao2024expel}, \ours{} stores high-level strategies and reasoning hints. By abstracting experiences into reusable reasoning units, \ours{} enables agents to generalize not only from successful cases~\citep{zhao2024expel, alazraki2025no, fu2024autoguide} but also by learning from failures, providing richer guidance for test-time learning. 

\noindent\textbf{Agent Test-Time Scaling.}
Test-time scaling (TTS)~\citep{snell2025scaling} has demonstrated strong effectiveness and has become a widely adopted practice in end-to-end problem-solving, such as coding~\citep{li2025s, yu2025z1} and math reasoning~\citep{muennighoff2025s1}, where methods including best-of-N~\citep{chow2025inferenceaware}, beam search~\citep{wu2024inference}, and leveraging verifiers~\citep{setlur2025scaling} are commonly employed. 
However, its application to multi-turn interactive scenarios, particularly agentic tasks, remains underexplored. Existing works mainly adapt the lessons learned from reasoning tasks~\citep{zhu2025scaling} and scale different dimensions of agentic systems, including the search space for each action~\citep{yu2025exact}, the number of agents in multi-agent systems~\citep{jin2025two}, and the number of interactions with the environment~\citep{shen2025thinking}. We found that none of these efforts considers the role of \textit{agent memory} in scaling, where an agent can learn from past experiences to guide future decisions. Our work extends this line of research by introducing memory-aware test-time scaling (\mats{}). As we will show in our empirical results (\S\ref{sec: res_mats} and \S\ref{sec: exp_synergy}), memory offers benefits beyond mere computational scaling, where memory and scaling synergistically work towards better performance.